\documentclass[a4paper,12pt]{article}
\usepackage{amsmath,amssymb,amsthm}
\usepackage[margin=1in]{geometry}

\newtheorem{theorem}{Theorem}
\newtheorem{lemma}[theorem]{Lemma}

\title{Problem 343: Fractional Sequences}
\author{Project Euler}
\date{}

\begin{document}
\maketitle

\section{Problem Statement}
For a positive integer $k$, define a sequence of fractions $a_i = x_i/y_i$ (in lowest terms) by $a_1 = 1/k$ and the iterative rule $a_{i+1} = (x_i + 1)/y_i$ reduced to lowest terms, terminating when $x_i = 0$. Let $f(k)$ be the last denominator before termination.

Compute $\displaystyle\sum_{\substack{p \text{ prime} \\ 5 \le p < 10^6}} f(p)$.

\section{Mathematical Foundation}

\begin{theorem}[Reduction to Fractional GCD]
The iterative fractional sequence is equivalent to a variant of the subtractive Euclidean algorithm. Each step reduces the pair $(x, y)$ via:
\[
(x, y) \;\longmapsto\; \bigl(y \bmod (x+1),\; x+1\bigr),
\]
terminating when the remainder is zero, at which point $f(k) = x + 1$.
\end{theorem}

\begin{proof}
At each step, the fraction $x/y$ is in lowest terms. The operation $(x+1)/y$ reduced to lowest terms advances the sequence. The standard division $y = q(x+1) + r$ with $0 \le r < x+1$ decomposes the step. When $r = 0$, we have $y \mid (x+1)$, so the fraction becomes an integer and the process terminates with last denominator $x+1$. Otherwise, the new pair $(r, x+1)$ has strictly smaller first component, guaranteeing convergence.
\end{proof}

\begin{lemma}[Logarithmic Convergence]
For input $k$, the sequence terminates in $O(\log k)$ steps.
\end{lemma}

\begin{proof}
Each step replaces $(x, y)$ with $(y \bmod (x+1), x+1)$ where $y \bmod (x+1) < x+1 \le y$. The quantity $\max(x,y)$ strictly decreases at each step, and the analysis mirrors that of the Euclidean algorithm: after two steps, the larger component decreases by at least a factor of $\phi = (1+\sqrt{5})/2$, giving $O(\log k)$ total steps.
\end{proof}

\section{Algorithm}

\begin{verbatim}
function f(k):
    x = 1, y = k
    while True:
        r = y mod (x + 1)
        if r == 0: return x + 1
        x, y = r, x + 1

function solve():
    for each prime p in [5, 10^6):
        total += f(p)
    return total
\end{verbatim}

\section{Complexity Analysis}
\begin{itemize}
    \item \textbf{Time:} $O\bigl(\pi(10^6) \cdot \log(10^6)\bigr) \approx O(1.6 \times 10^6)$, plus $O(N \log\log N)$ for sieving.
    \item \textbf{Space:} $O(N)$ for the sieve, $N = 10^6$.
\end{itemize}

\section{Answer}

\[
\boxed{269533451410884183}
\]

\end{document}
